% ---
% id: statistics/def:taxonomy
% kind: def
% title: "Admissible geometries and the canonical box form"
% summary: >-
%   An admissible geometry is a singleton, one rational vector, or a polytope,
%   rational linear inequalities intersected with the probability simplex, each
%   with its validation predicate. A polytope in canonical box form is given by
%   recorded lower and upper bounds per outcome.
% prominence: standard
% status: implemented
% uses: []
% assumes: []
% refs: []
% provenance: newly-derived
% implements: certiqs_key.statistics.taxonomy:GeometryTag
% lean:
% ---

The acceptance test asks whether the frequency vector of a run lies in a
declared set. The set is declared in one of two forms, and a polytope may in
addition record per-outcome bounds from which its rows are built.

\begin{definition}\label{def:taxonomy}
An \term[A tag, singleton or polytope, together with the data of that tag; an
acceptance set is declared as one.]{admissible geometry} is a pair consisting of a \emph{tag}
and its \emph{data}, drawn from exactly the following two cases, each with
its data, the set $\Gamma$ the data denote and its
\term[The condition on a geometry's data that validation checks: a
probability vector for a singleton, a feasible linear system for a
polytope.]{validation predicate}. There is no third tag.
\begin{description}[leftmargin=1.9em]
\item[$\mathrm{singleton}$.] Data: a vector
$\mathbf{q}^{\ast}\in\mathbb{Q}^{\Chatset}$. Denotation
$\Gamma\coloneqq\{\mathbf{q}^{\ast}\}$. \emph{Validation predicate:}
$q^{\ast}(\hat c)\ge0$ for every $\hat c$ and $\sum_{\hat c}q^{\ast}(\hat
c)=1$.
\item[$\mathrm{polytope}$.] Data: an integer row count $m_{\mathrm{row}}\ge0$,
a matrix $B\in\mathbb{Q}^{\,m_{\mathrm{row}}\times\Chatset}$ and a vector
$\mathbf{b}\in\mathbb{Q}^{\,m_{\mathrm{row}}}$. Denotation
\begin{equation}\label{eq:polytope}
  \Gamma\coloneqq\{\mathbf{q}\in\simplex{\Chatset}:\ B\mathbf{q}\le\mathbf{b}\},
\end{equation}
the inequality componentwise. \emph{Validation predicate:} feasibility of the
linear system $\{\mathbf{q}\ge0,\ \sum_{\hat c}q(\hat c)=1,\
B\mathbf{q}\le\mathbf{b}\}$ over the rationals.
\end{description}
A $\mathrm{polytope}$ geometry is in \term[A polytope geometry that records a
lower and an upper bound per outcome and whose rows are exactly the two
bounds of each outcome.]{canonical box form} when it carries, in addition, a
pair $\qlo,\qhi\in\mathbb{Q}^{\Chatset}$ with $0\le\qlo\le\qhi$
componentwise, and its $(B,\mathbf{b})$ are exactly
\begin{equation}\label{eq:boxrows}
  m_{\mathrm{row}}=2m,\qquad
  \begin{aligned}
    &\text{row }i:\ \ \delta_{\hat c^{(i)}}^{\transpose},
      &&\text{right-hand side } q^{+}(\hat c^{(i)}),\\
    &\text{row }m+i:\ \ -\delta_{\hat c^{(i)}}^{\transpose},
      &&\text{right-hand side } -q^{-}(\hat c^{(i)}),
  \end{aligned}
  \qquad i=1,\dots,m,
\end{equation}
where $m\coloneqq|\Chatset|$, $\hat c^{(1)},\dots,\hat c^{(m)}$ is the
declaration's enumeration of $\Chatset$ and $\delta_{\hat c}$ is the standard
basis vector of $\mathbb{R}^{\Chatset}$ at $\hat c$. Row indices are integers,
never alphabet elements; equivalently
\begin{equation}\label{eq:boxset}
  \Gamma=\{\mathbf{q}\in\simplex{\Chatset}:\
  q^{-}(\hat c)\le q(\hat c)\le q^{+}(\hat c)\ \ \forall\hat c\in\Chatset\}.
\end{equation}
\end{definition}
